\chapter{A benchmark CSV fájlok formátuma}

A projektben minden mérés egy külön CSV fájlba kerül. A fájl két részből áll.

\section{Meta szekció}

A meta sorok \code{\#Kulcs,Érték} formában jelennek meg, például:

\begin{verbatim}
#Algorithm,Gcm
#KeySizeBits,256
#DataSizeMegabytes,64
#FrameworkDescription,.NET 10.0.5
#WindowsBuild,Build 26200.8037
#BuildArchitecture,x64
#PowerSource,AC power
#OpenClDeviceName,NVIDIA GeForce RTX 4060
\end{verbatim}

Ezeket a sorokat a LaTeX oldalon egyszerűen lehet ignorálni
(\code{comment char=\#}), illetve a jelen dokumentum a metaadatokból
külön összefoglaló táblát is készít (\code{generated/machines.csv} és a tömörebb \code{generated/machines_summary.csv}).

\section{Adatszekció}

A meta után jön a tényleges táblázat, például:

\begin{verbatim}
TimestampUtc,Iteration,Engine,Direction,Algorithm,Padding,KeySizeBits,...
2026-03-18T19:37:59Z,1,NativeCpu,Encrypt,Gcm,None,256,...,29.94,True
2026-03-18T19:37:59Z,1,OpenCl,Encrypt,Gcm,None,256,...,690.58,True,23.06
\end{verbatim}

A dokumentum készítésekor nem közvetlenül a nyers sorokat rajzolom ki,
hanem iterációk szerint átlagot/szórást számolok, mert így a grafikonok és táblák
áttekinthetőek és kevésbé zajosak. A nyers adatok természetesen megmaradnak.
